% =========================================================
% Breadcrumb
% =========================================================
\begin{tcolorbox}[
  colback=gray!6,
  colframe=gray!40,
  arc=2pt,
  left=8pt, right=8pt, top=6pt, bottom=6pt,
  title={\small\textbf{Where You Are in the Journey}},
  fonttitle=\small\bfseries
]
\begin{center}
\small
Propositional Logic
$\;\to\;$ Sets \& Functions
$\;\to\;$ $\mathbb{R}$
$\;\to\;$ Algebraic Structures
$\;\to\;$ \textbf{Set Algebras}
$\;\to\;$ Measure Theory
$\;\to\;$ $\cdots$
\end{center}

\medskip
\noindent\textbf{How we got here.}
Set theory gave us the basic operations --- union, intersection, complement,
and power set --- and the language of subsets of a fixed space $X$.
Algebraic structures showed us how to classify mathematical objects by
their closure properties under operations. Set algebras bring these two
threads together: we ask which \emph{collections} of subsets are closed
under which set-theoretic operations, and what algebraic laws govern them.

\medskip
\noindent\textbf{What this chapter builds.}
We develop the hierarchy of closure systems on $2^X$: rings of sets,
algebras (fields) of sets, $\sigma$-rings, and $\sigma$-algebras.
Alongside this, we study characteristic functions as a bridge between
subsets and arithmetic, revealing that $2^X$ carries Boolean ring and
$\mathbb{F}_2$-vector space structure. The power set itself is the
ambient $\sigma$-algebra and the universal example.

\medskip
\noindent\textbf{Where this leads.}
Measure theory requires a $\sigma$-algebra as its foundational input:
a measure is defined on a $\sigma$-algebra, and measurability is a
membership condition in one. The Boolean algebra structure of $2^X$
connects back to propositional logic. Characteristic functions
reappear as simple functions in Lebesgue integration theory.
\end{tcolorbox}
\vspace{1em}

% =========================================================
% Structural Roadmap
% =========================================================
\section*{Structural Roadmap}

Each major topic is organised as:
\begin{center}
\textbf{Definitions $\longrightarrow$ Main Theorems
$\longrightarrow$ Consequences and Structural Insight}
\end{center}

The global progression is:
\begin{enumerate}
  \item Families of sets and closure operations
  \item Finite closure structures: rings of sets and algebras of sets
  \item Countable closure structures: $\sigma$-rings and $\sigma$-algebras
  \item The power set as a Boolean ring and $\mathbb{F}_2$-vector space
  \item Characteristic functions and the identification $2^X \cong \{0,1\}^X$
\end{enumerate}

\begin{remark}[Primary sources]
The treatment follows Kolmogorov and Fomin \textit{Introductory Real Analysis},
Chapter~1, and standard accounts in measure theory (Rudin, Royden).
Characteristic functions and the Boolean ring structure appear in Halmos
\textit{Measure Theory}.
\end{remark}

% =========================================================
% Content
% =========================================================
\input{volume-iv/algebra/set-algebras/notes/notes-set-algebras}
% =========================================================
% The Power Set and Characteristic Functions
% =========================================================

\section{The Power Set and Characteristic Functions}


% Future sections (uncomment as content expands):
% \input{volume-iv/algebra/set-algebras/notes/notes-generated-systems}
% \input{volume-iv/algebra/set-algebras/notes/notes-auxiliary-systems}